% Bagian: first-order-logic
% Bab: models-theories
% Subbagian: expressing-props-of-structures

\documentclass[../../../include/open-logic-section]{subfiles}

\begin{document}

\olfileid[id]{fol}{mat}{exs}

\olsection{Mengungkapkan Sifat-Sifat \printtoken{P}{structure}}

\begin{explain}
Sering kali berguna dan penting untuk mengungkapkan syarat-syarat pada fungsi
dan relasi, atau secara lebih umum, bahwa fungsi dan relasi dalam suatu
struktur memenuhi syarat-syarat tersebut. Misalnya, kita ingin mempunyai cara
untuk membedakan !!{structure}s bagi suatu bahasa yang ``menangkap'' makna
yang kita kehendaki bagi !!{predicate}s dari struktur yang tidak demikian.
Tentu saja kita sepenuhnya bebas menentukan !!{structure}s mana yang kita
``maksudkan''; misalnya, kita dapat menentukan bahwa interpretasi
!!{predicate}~$\le$ harus berupa suatu urutan, atau bahwa kita hanya tertarik
pada interpretasi~$\Lang L$ yang domainnya terdiri atas himpunan dan yang
menafsirkan $\Obj \in$ sebagai relasi ``merupakan !!a{element} dari.'' Namun,
dapatkah kita melakukan ini dengan !!{sentence}s bahasa tersebut? Dengan kata
lain, syarat-syarat mana pada !!a{structure}~$\Struct M$ yang dapat kita
ungkapkan dengan !!a{sentence} (atau mungkin suatu himpunan !!{sentence}s)
dalam bahasa~$\Struct M$? Ada beberapa syarat yang tidak akan dapat kita
ungkapkan. Misalnya, tidak ada kalimat dari~$\Lang L_A$ yang hanya benar dalam
!!a{structure}~$\Struct M$ jika $\Domain M = \Nat$. Kita tidak dapat
menyatakan ``domain hanya memuat bilangan asli.'' Namun, mungkin ada
``sifat-sifat struktural'' dari !!{structure}s yang dapat kita ungkapkan.
Sifat-sifat !!{structure}s mana yang dapat kita ungkapkan dengan
!!{sentence}s? Atau, dengan rumusan lain, kumpulan !!{structure}s mana yang
dapat kita gambarkan sebagai semua struktur yang membuat !!a{sentence} (atau
suatu himpunan !!{sentence}s) benar?
\end{explain}

\begin{defn}[Model suatu himpunan]
Misalkan $\Gamma$ suatu himpunan !!{sentence}s dalam bahasa~$\Lang L$. Kita
mengatakan bahwa !!a{structure}~$\Struct M$ \emph{merupakan model
dari}~$\Gamma$ jika $\Sat{M}{!A}$ bagi semua $!A \in \Gamma$.
\end{defn}

\begin{ex}
Kalimat $\lforall[x][x \le x]$ benar dalam~$\Struct M$ jika dan hanya jika
$\Assign{\le}{M}$ merupakan relasi refleksif. Kalimat
$\lforall[x][\lforall[y][((x \le y \land y \le x) \lif x = y)]]$ benar
dalam~$\Struct M$ jika dan hanya jika $\Assign{\le}{M}$ antisimetris.
Kalimat $\lforall[x][\lforall[y][\lforall[z][((x \le y \land y \le z)
      \lif x \le z)]]]$ benar dalam~$\Struct M$ jika dan hanya jika
$\Assign{\le}{M}$ transitif. Dengan demikian, model-model dari
\begin{align*}
\{\quad &\lforall[x][x \le x], \\
   & \lforall[x][\lforall[y][((x \le y \land y \le
    x) \lif x = y)]], \\
   &\lforall[x][\lforall[y][\lforall[z][((x \le y
      \land y \le z) \lif x \le z)]]] \quad \}
\end{align*}
tepat merupakan struktur-struktur yang di dalamnya~$\Assign{\le}{M}$
refleksif, antisimetris, dan transitif, yakni suatu urutan parsial. Oleh karena
itu, kita dapat mengambil kalimat-kalimat tersebut sebagai aksioma bagi
\emph{teori orde pertama tentang urutan parsial}.
\end{ex}

\end{document}
